\chapter*{Appendix: Recommended Resources}
\addcontentsline{toc}{chapter}{Appendix: Recommended Resources}

The following books, researchers, and organizations have informed or are 
closely aligned with the principles in this book. Each offers a valuable 
next step for the reader who wants to go deeper.

\section*{Books}

\begin{description}[leftmargin=2em, style=nextline]

  \item[Daniel J.\ Siegel \& Mary Hartzell, \textit{Parenting from the 
    Inside Out} (2003)]
    A compelling exploration of how our own childhood experiences shape 
    our parenting---and how to parent more consciously.

  \item[Daniel J.\ Siegel \& Tina Payne Bryson, \textit{The Whole-Brain 
    Child} (2011)]
    A practical and accessible guide to integrating the left and right 
    hemispheres of the developing brain through everyday interactions.

  \item[John Gottman \& Joan DeClaire, \textit{Raising an Emotionally 
    Intelligent Child} (1997)]
    The definitive guide to emotion coaching, grounded in Gottman's 
    longitudinal research.

  \item[Marshall B.\ Rosenberg, \textit{Nonviolent Communication: A 
    Language of Life} (2003)]
    The original and essential text on NVC---practical, provocative, 
    and life-changing.

  \item[Ross W.\ Greene, \textit{The Explosive Child} (1998)]
    A revolutionary approach to understanding and helping children who 
    seem chronically challenging, based on collaborative problem-solving.

  \item[Lawrence Cohen, \textit{Playful Parenting} (2001)]
    A joyful argument for the therapeutic and relational power of play 
    in family life.

  \item[Bren\'e Brown, \textit{Daring Greatly} (2012)]
    While not exclusively a parenting book, Brown's research on 
    vulnerability, shame, and connection is foundational to empathetic 
    parenting.

  \item[Alfie Kohn, \textit{Unconditional Parenting} (2005)]
    A thought-provoking challenge to conventional reward-and-punishment 
    approaches and a case for love without strings.

  \item[Carol S.\ Dweck, \textit{Mindset: The New Psychology of Success} 
    (2006)]
    Essential reading on the growth mindset and how parents can foster 
    it in children.

  \item[Adele Faber \& Elaine Mazlish, \textit{How to Talk So Kids Will 
    Listen \& Listen So Kids Will Talk} (1980)]
    A classic---still among the most practical and widely used 
    communication guides for parents.


  \item[Jean Decety \& C\'eline Holvoet, \textit{The Neurobiology of Empathy in Infancy} (2021)]
    An exploration of the neural circuits responsible for empathy, demonstrating they are active early in life.

  \item[Giacomo Rizzolatti \& Laila Craighero, \textit{The Mirror-Neuron System} (2004)]
    The foundational paper on the biological mechanism of mirror neurons and how they bridge the gap between self and other.

  \item[Felix Warneken, \textit{How Children Solve the Two Challenges of Cooperation} (2018)]
    Compelling evidence of intrinsic altruism and prosociality in toddlers without the expectation of reward.

  \item[Martin L.\ Hoffman, \textit{Empathy and Moral Development: Implications for Caring and Justice} (2000)]
    A comprehensive model of social-emotional development, outlining the stages of empathy from infancy to maturity.

  \item[Ruth Feldman, \textit{The Neurobiology of Human Attachments} (2024)]
    Recent advancements in Attachment Theory, highlighting bio-behavioral synchrony and interbrain networks.

\end{description}

\section*{Key Researchers and Theorists}

\begin{description}[leftmargin=2em, style=nextline]
  \item[John Bowlby \& Mary Ainsworth] Founders of attachment theory.
  \item[Erik Erikson] Stages of psychosocial development.
  \item[Alexander Thomas \& Stella Chess] Temperament research.
  \item[Daniel Goleman] Popularizer of emotional intelligence.
  \item[Peter Salovey \& John Mayer] Original EI researchers.
  \item[Dan Siegel] Interpersonal neurobiology and parenting.
  \item[John Gottman] Emotion coaching and family research.
\end{description}

\section*{Organizations and Online Resources}

\begin{itemize}
  \item \textbf{The Gottman Institute} --- \url{https://www.gottman.com}
  \item \textbf{Center for Nonviolent Communication} --- \url{https://www.cnvc.org}
  \item \textbf{Circle of Security International} --- 
        \url{https://www.circleofsecurityinternational.com}
  \item \textbf{Zero to Three} --- \url{https://www.zerotothree.org}
  \item \textbf{Collaborative for Academic, Social, and Emotional 
        Learning (CASEL)} --- \url{https://casel.org}
\end{itemize}
